\section{Introdução}
\label{sec:introducao}

A computação em nuvem consolidou-se como uma das principais estratégias adotadas pelas organizações para modernizar infraestrutura, ganhar escala e otimizar recursos computacionais. Nesse cenário, a cultura \textit{DevOps} ganhou destaque ao integrar desenvolvimento e operações e incentivar automação, integração contínua e entrega contínua \cite{kim2016devops}, e a \textit{Infrastructure as Code} (\textit{IaC}) se encaixa nesse contexto ao permitir que recursos de infraestrutura sejam definidos e gerenciados por código versionado, reutilizável e passível de automação \cite{morris2016infrastructure}. Paralelamente, o avanço de ferramentas de Inteligência Artificial generativa, capazes de transformar descrições em linguagem natural em código \cite{chen2021codex}, passou a ser explorado também na criação de artefatos de infraestrutura, como alternativa para reduzir o esforço técnico exigido pela \textit{IaC}.

Apesar desse potencial, reduzir o esforço de desenvolvimento não garante que os \textit{templates} produzidos por IA atendam aos requisitos de padronização, conformidade e segurança esperados no desenvolvimento manual, sobretudo quando as regras internas da organização não são fornecidas ao modelo durante a geração. \textbf{O problema investigado neste trabalho é, portanto, compreender como o uso de Inteligência Artificial impacta a eficiência operacional, a padronização e a segurança de \textit{templates} de \textit{IaC}, quando comparado ao desenvolvimento manual e à geração com contexto organizacional.}

Essa questão é relevante porque incorporar novos recursos de nuvem não pode comprometer requisitos técnicos e de segurança. Artefatos de infraestrutura podem apresentar permissões excessivas, exposição de dados e outros problemas de configuração \cite{rahman2021securitysmells,verdet2025security}, além de defeitos de reconciliação do estado da infraestrutura \cite{hassan2024state}, e processos dependentes de atividades manuais podem aumentar o tempo das entregas \cite{forsgren2018accelerate}, o que torna relevante avaliar alternativas que reduzam o esforço de criação de \textit{templates} sem comprometer sua qualidade.

\textbf{O objetivo geral deste trabalho é avaliar o uso da Inteligência Artificial na geração de \textit{templates} de Infraestrutura como Código, considerando eficiência operacional, padronização e segurança.} Para alcançá-lo, foram definidos três objetivos específicos: (i) analisar se a IA reduz o tempo e o esforço necessários para a criação dos artefatos, (ii) avaliar a conformidade dos \textit{templates} gerados com boas práticas e padrões de \textit{IaC}, e (iii) investigar vulnerabilidades de segurança produzidas pelas diferentes abordagens de geração.

Com essa avaliação, buscou-se produzir evidências empíricas sobre os benefícios e as limitações do uso de Inteligência Artificial na geração de artefatos de infraestrutura, em particular sobre até que ponto o fornecimento de contexto organizacional aproxima os \textit{templates} gerados por IA dos padrões de qualidade obtidos no desenvolvimento manual. Essas evidências contribuem para a adoção mais segura e eficiente dessas tecnologias em ambientes \textit{DevOps}.

O restante deste artigo está organizado da seguinte forma. A Seção~\ref{sec:fundamentacao} apresenta os conceitos de \textit{DevOps}, \textit{IaC}, Inteligência Artificial generativa e \textit{CI/CD} que fundamentam o trabalho. A Seção~\ref{sec:trabalhos_relacionados} discute estudos relacionados à qualidade, segurança, testes e geração automatizada de artefatos \textit{IaC}. A Seção~\ref{sec:metodologia} descreve a metodologia adotada para geração, validação e comparação dos \textit{templates}, a Seção~\ref{sec:resultados} apresenta os resultados obtidos, a Seção~\ref{sec:ameacas} discute as ameaças à validade e a Seção~\ref{sec:conclusao} apresenta as considerações finais e os trabalhos futuros.
